\begin{helper}[A repeated color on a path-shaped minimum bad tree]
Let \((\varphi,T)\) be a minimum bad pair with
\(T=(y_0,e_1,y_1,\ldots,e_n,y_n)\), \(n\geq2\), and \(\rho(T)=1\).
Then there are a color \(\beta\) and an index \(i<n\) such that
\(\beta\) is missing at both \(y_i\) and \(y_n\).
\end{helper}
\begin{proof}
Let \(P=T_{n-1}\) be the object \noderef{TreeSequencePrefix} obtained by
retaining exactly the vertices \(y_0,\ldots,y_{n-1}\) and edges
\(e_1,\ldots,e_{n-1}\).  Since \(n\geq2\), this is a positive proper
initial segment of \(T\).  Therefore
\noderef{TashTreeProperPrefixElementary}, applied to the supplied
\noderef{TashTreeMinimumBadPair}, says that
\noderef{TreeSequenceVertexSet} of \(P\) is elementary under
\(\varphi\).

On the other hand, unfolding the badness conjunct with
\noderef{TashTreeBadPair} says that
\noderef{TreeSequenceVertexSet} of \(T\) is not elementary.  By the
definition \noderef{ColoredVertexSetElementary}, there are distinct
vertices \(u,v\) of \(T\) and a color \(\beta\) which lies in both of
their \noderef{MissingColors} sets.  If both vertices belonged to \(P\),
the same pair and color would contradict the elementariness of \(P\).
Thus one of \(u,v\) is the only displayed vertex outside \(P\), namely
\(y_n\).  The other is a displayed vertex \(y_i\) with \(i<n\), because
\noderef{TreeSequenceVertexSet} is precisely the image of the displayed
vertex positions.  Relabeling the two witnesses if necessary gives the
asserted \(\beta,i\).  (The path-index hypothesis records the branch in
which this witness will be used; the extraction itself needs only the
minimum bad pair and its proper-prefix elementariness.)
\end{proof}
